\documentclass{article}
\usepackage{tabularx}
\newcolumntype{b}{X}
\newcolumntype{s}{>{\hsize=.4\hsize}X}
\title{Return to Farlorret's ruling on the interaction between the Fighting style feat: Two-weapon fighting \& the General feat: Dual wielder}
\date{\today}



\begin{document}
    \maketitle

    \section*{Background}

    \textbf{Two-Weapon Fighting}\\
    Fighting Style Feat (Prerequisite: Fighting Style Feature)

    When you make an extra attack as a result of using a weapon that has the Light property, you can add your ability modifier to the damage of that attack if you aren’t already adding it to the damage.\\
    \textbf{Dual Wielder}\\
    General Feat (Prerequisite: Level 4+, Strength or Dexterity 13+)

    You gain the following benefits.
    
    \textbf{Ability Score Increase.} Increase your Strength or Dexterity score by 1, to a maximum of 20.
    
    \textbf{Enhanced Dual Wielding.} When you take the Attack action on your turn and attack with a weapon that has the Light property, you can make one extra attack as a Bonus Action later on the same turn with a different weapon, which must be a Melee weapon that lacks the Two-Handed property. You don’t add your ability modifier to the extra attack’s damage unless that modifier is negative.
    
    \textbf{Quick Draw.} You can draw or stow two weapons that lack the Two-Handed property when you would normally be able to draw or stow only one.\\
    \textbf{Nick}

    When you make the extra attack of the Light property, you can make it as part of the Attack action instead of as a Bonus Action. You can make this extra attack only once per turn.\\
    \textbf{Light}

    When you take the Attack action on your turn and attack with a Light weapon, you can make one extra attack as a Bonus Action later on the same turn. That extra attack must be made with a different Light weapon, and you don’t add your ability modifier to the extra attack’s damage unless that modifier is negative. For example, you can attack with a Shortsword in one hand and a Dagger in the other using the Attack action and a Bonus Action, but you don’t add your Strength or Dexterity modifier to the damage roll of the Bonus Action unless that modifier is negative.
    \section*{Ruling}

    The added ability modifier from Two-Weapon Fighting is only applicable on attacks made due to the interaction stated in the light property.

    This excludes the case where the Dual Wielder feat allows for an extra attack to be made as a bonus action from benefitting from Two-Weapon Fighting.

    Additionally, any possible interactions where this Fighting style could be interpreted to apply outside of the light property due to combinations of weapons, properties or features does not.
    \section*{Examples}
    \begin{table}[h]
        \begin{tabularx}{\textwidth}{|sss|bb|}
            \hline
            Nick & Two-Weapon Fighting & Dual Wielder & Action & Bonus Action\\\hline
            No&No&No&One attack with a Light weapon with Damage modifier& One attack with a different Light weapon without Damage modifier\\\hline
            Yes&No&No&Two attacks with two different Light weapons only one with Damage modifier.&\\\hline
            No&Yes&No&One attack with a Light weapon with Damage modifier& One attack with a different Light weapon with Damage modifier\\\hline
            No&No&Yes&One attack with a Light weapon with Damage modifier& One attack with a different Light weapon without Damage modifier\\\hline
            Yes&Yes&No&Two attacks with two different Light weapons Both with Damage modifier.&\\\hline
            Yes&No&Yes&Two attacks with two different Light weapons only one with Damage modifier.&One attack with a different Light weapon without Damage modifier\\\hline
            No&Yes&Yes&One attack with a Light weapon with Damage modifier& One attack with a different Light weapon with Damage modifier\\\hline
            Yes&Yes&Yes&Two attacks with two different Light weapons Both with Damage modifier.&One attack with a different Light weapon without Damage modifier\\\hline
            
        \end{tabularx}
    \end{table}
    \end{document}